%!TEX root = ../template.tex
%%%%%%%%%%%%%%%%%%%%%%%%%%%%%%%%%%%%%%%%%%%%%%%%%%%%%%%%%%%%%%%%%%%%
%% chapter3-background-draft.tex
%%
%% DRAFT: starting point of this work — which components are inherited and
%% which are contributed here. Intended to open Chapter 3, before
%% Section~\ref{sec:victim_architectures}.
%%%%%%%%%%%%%%%%%%%%%%%%%%%%%%%%%%%%%%%%%%%%%%%%%%%%%%%%%%%%%%%%%%%%

% ─────────────────────────────────────────────────────────────────────────────
\section{Starting Point and Scope of This Work}
\label{sec:starting_point}

This work is the third stage of a line of research, and it is useful to state at
the outset which components are inherited and which are contributed here, since the
research question is comparative and therefore depends on the earlier stages
remaining unchanged.

\subsection{The routing system under test}
\label{subsec:base_system}

The victim is a multi-agent routing controller developed in earlier work within the
same research line~\cite{amaral2021maddpg-routing}. It applies Multi-Agent Deep
Deterministic Policy Gradient (MADDPG) to adaptive packet routing under the
paradigm of centralised training with decentralised execution
(CTDE)~\cite{lowe2017multi}: each agent $i$ holds an actor $\pi_{\theta_i}$ that
conditions only on its own local observation, while its critic $Q_{\phi^i}$ may,
during training, observe the joint state and the actions of every other agent.
Conditioning the critic on the joint action is what stabilises learning, since from
any single agent's perspective the environment would otherwise appear
non-stationary as the other agents' policies evolve~\cite{lowe2017multi}. After
training the critics are discarded and each agent acts on local information alone.
In this routing formulation every forwarding node is an agent whose observation
summarises local link utilisation, queue occupancy and the destinations awaiting
service, and whose action selects, for each destination, one of $K$ pre-computed
candidate paths; the trained system is thus a distributed controller in which no
node holds a global view, yet the paths chosen collectively determine end-to-end
delivery. That prior work also investigated critic design and graph-based state
encoding, which is the origin of the variant family evaluated in
Section~\ref{sec:victim_architectures}.

This thesis does not modify the routing controller. The trained agents are treated
as a fixed system under test: their parameters are frozen, and the adversary
interacts with them only through forward evaluations and gradients taken with
respect to their inputs.

\subsection{The FGSM baseline}
\label{subsec:base_fgsm}

The second stage established the adversarial setting. A single-step attack based on
the Fast Gradient Sign Method
(FGSM)~\cite{goodfellow2015explainingharnessingadversarialexamples} was developed
against these agents in a preceding dissertation in the same line of
work~\cite{TODO-colleague-fgsm}. FGSM perturbs an input one step along the sign of
the gradient of a loss with respect to that input,
$\eta = \varepsilon \cdot \operatorname{sign}\!\big(\nabla_x J(\theta, x, y)\big)$,
under an $\ell_\infty$ budget $\lVert \eta \rVert_\infty \le \varepsilon$. Applied
to the routing controller, the victim is the trained actor, the input is an agent's
local observation, and the objective rewards actions that degrade routing
performance rather than ones that induce a misclassification. The threat model it
fixed is \emph{observation-space} rather than actuation-space: the adversary
corrupts what the agents perceive, modelling a compromised telemetry or monitoring
plane, while the network itself continues to evolve according to its true state.

That stage contributed not only an attack but an evaluation protocol --- the
perturbation budget sweep, the frozen-victim assumption, and the metrics by which
an attack is judged --- and it returned a negative result. The FGSM perturbation
alters a substantial fraction of the routing decisions taken by the victim, yet
recovers only a small part of the degradation achievable in principle, because the
redundancy of the $K$-path topology absorbs the redirected traffic and leaves
end-to-end delivery largely intact.

\subsection{Contribution of this thesis}
\label{subsec:contribution_scope}

That negative result admits two readings: either the controller is genuinely
robust, or a single-step, per-agent adversary is too weak for its failure to
constitute evidence of robustness. Distinguishing between them requires a stronger
adversary, and supplying one is the contribution of this work. Three extensions are
developed, all within the inherited threat model: replacing the single gradient
step with an iterated projected-gradient
optimisation~\cite{madry2019deeplearningmodelsresistant}; coordinating the
perturbation jointly across compromised agents rather than treating each in
isolation (Section~\ref{subsec:coordinated_attack}); and concentrating a limited
attack budget on high-leverage moments rather than spending it uniformly
(Section~\ref{subsec:timed_attack}).

Because the purpose is comparison rather than a stronger attack in isolation,
compatibility with the earlier stage is a design constraint rather than a
convenience. The attacks developed here operate on the same frozen agents, within
the same $\ell_\infty$ ball and the same domain constraints, and are scored under
the same evaluation protocol as the FGSM baseline. Any difference in measured
damage is then attributable to the attack algorithm itself, rather than to a
change in the adversary's capabilities or in the conditions of measurement; where
this constraint forced a specific implementation choice, it is noted at the
relevant point in the text.
